\section{Support-First Scheduling and Flexible Capacity}

A small support edit enables a large gain edit on the same
or a later query. The retrospective timing control retains the
same support action and role counts in each recording, but
uses the earliest positive-gain row distinct from that support
recipient. The direct lower contrast is \ClockSupportEffect{}.
The online rule makes \ClockSupportAnchorEdits{} edits and
passes \ClockSupportAnchorJoint{} views, versus
\ClockSupportMatchedJoint{} with the same budget. The old
independent schedule passes \ClockSupportIndependentJoint{}
at \ClockSupportIndependentEdits{} edits, a different budget.
A fixed learned support threshold spends only
\ClockSupportLearnedEdits{} edits and passes
\ClockSupportLearnedJoint{} views. None meets the original risk rule.

Allowing both roles to change time and recipient at each
recording's original role counts gives \ClockSupportAny{}/606
verified outcome-aware witnesses; imposing support before or
with gain leaves \ClockSupportOrdered{}. With one edit per role
and flexible realized counts, the limit is \ClockSupportAtMost{}.
There are \ClockSupportNoReference{} views with no full easy
reference. Another \ClockSupportOLSBlocked{} have a risk-feasible
improvement over CV but none over OLS4 in the neural bank.
Adding OLS4 and observed-prefix directions at equal per-row
peak displacement, without amplification or new eligible rows,
raises capacity to \ClockSupportShape{}. These bounds diagnose
the bank and schedule, not causal performance.
